% Source: https://artofproblemsolving.com/community/c3h1592098_mathematics_question_from_brilliant
% Stable identity: aops-topic-1592098-q01
There are $\dbinom{14}{4}=1001$ equally likely choices. Count the complement: choices with no two consecutive numbers.

\textbf{A bijection.} Write a choice with no two consecutive numbers as $a<b<c<d$, so $b\ge a+2$, $c\ge b+2$, $d\ge c+2$. Then
\[
a<b-1<c-2<d-3
\]
are four distinct numbers from $\{1,\dots,11\}$. Conversely, from any four numbers $a'<b'<c'<d'$ in $\{1,\dots,11\}$ we get back $a',\,b'+1,\,c'+2,\,d'+3$, four numbers from $\{1,\dots,14\}$ with no two consecutive. So there are $\dbinom{11}{4}=330$ such choices.

Hence
\[
P(\text{some two consecutive})=1-\dfrac{330}{1001}=\dfrac{671}{1001}=\boxed{\dfrac{61}{91}}.
\]
